we define a @@M0@@-homomorphism: @@M1@@ extending the augmentation of @@M2@@ by sending @@M3@@ to @@M4@@, which endows @@M5@@ with an augmented @@M6@@-@@M7@@-algebra structure. We propose to determine its filtration (3.10) in terms of that of @@M8@@.

We grade @@M9@@ by weight with respect to the @@M10@@, assigning @@M11@@ the weight @@M12@@; we denote by @@M13@@ the set of monomials of weight @@M14@@ in the indeterminates @@M15@@.

\sgaStableTarget{sga6.e05.lem.4-11}\textbf{Lemma 4.11.} \uline{Let} @@M16@@; \uline{the coefficient of a monomial of weight} @@M17@@ \uline{with respect} \uline{to the} @@M18@@ \uline{in} @@M19@@ \uline{has filtration at least} @@M20@@ (\uline{in} @@M21@@)

The polynomial @@M22@@ has no isobaric components of weight @@M23@@ in the @@M24@@, @@M25@@ (3.5), so the lemma is true when @@M26@@ has zero augmentation.

If @@M27@@ denotes the augmentation homomorphism: @@M28@@, we can write in the general case:
@@M29@@
Now @@M30@@, the @@M31@@ being isobaric of weight @@M32@@ in the @@M33@@ (3.8); we deduce the lemma by taking the product.

\sgaStableTarget{sga6.e05.prop.4-12}\textbf{Proposition 4.12.} \sgaStableTarget{sga6.e05.eq.4-12-1}@@M34@@.

It is clear that the second ideal is contained in the first; it therefore remains only to prove the reverse inclusion.

Since the filtration defined by the family of ideals on the right-hand side for varying @@M35@@ is a ring filtration, it suffices to prove the reverse inclusion for an element of the form @@M36@@, with @@M37@@, these elements
